\section{项目定位、范围与当前结论}
\label{sec:scope}
\subsection{一项审查判断究竟回答什么}
cv-agent 面向明确授权的本地源码审查。最小判断对象不是一个孤立的 sink 字符串，而是源码快照、候选位置、待判性质、前提及范围的组合：
\begin{equation}
J=(R,c,p,\Gamma,\Omega),
\end{equation}
其中 $R$ 是源内容身份，$c$ 是候选及其原始行号，$p$ 是机器可核对的性质，$\Gamma$ 是明确输入、部署和环境前提，$\Omega$ 是覆盖范围。同一函数中的不同调用、同一源码下的不同前提、权限 mode 与完整授权边界，分别是不同判断对象。任何证明都只能解释自己实际覆盖的对象。\cite{contract}

\subsection{当前存在的两条路径}
\textbf{产品路径}是新建的 \code{product_triage_v1}。它从使用者选择的 Python 源码目录出发，进行可读性/能力预检、字节清单与 AST 索引、确定性候选发现、适用静态检查、证据归并和产物生成。它不依赖 API key，也不执行目标模块。即使某个性质没有受支持的验证语义，仍能输出源码位置、局部观察、缺项和人工复核建议。

\textbf{研究路径}保留 E1--E5、可选 Jev、模型与工具循环、注册夹具和正向有界见证。它研究模型如何选择取证动作，以及检索、规划、多专家与提前终止是否带来净收益。其旧配置、预算、历史账本和 G-A 判定未被产品路径重新解释。共享源码完整性缺陷已在原边界修复，历史实验结果不会因此自动成为当前代码成绩。

两条路径的共同原则是：先确定判断对象；区分线索、事实、证明与标签；相关未知不得解释为安全；预算不改变证明门槛；工程完成与科学结论分别验收。

\subsection{本文的贡献与证据强度}
本文的核心结论是：\textbf{当前项目已具备独立验收的窄静态交付路径，同时原模型自主确认与外部迁移目标仍未成立。}支撑这项结论的证据包括：已复现的源码缺陷及正反控制、最终完整离线套件、真实 CLI 产物、拒绝式传输监测，以及保留的历史模型账本。它们各自回答不同问题。

\begin{table}[htbp]\centering
\begin{tabularx}{\textwidth}{L{29mm}L{45mm}Y}
\toprule 证据类别 & 本文依据 & 可支持的结论\\\midrule
源码实现 & 修后未提交工作树、锁定依赖 & 哪个入口、状态、算法和边界已实现。\\
离线回归 & 2684 passed；启用能力组 126 passed & 已枚举输入与故障控制通过，不是语义完备性。\\
CLI 验收 & 七个自有无害用例；额外 transport spy & 输入到合法产物闭环可工作，不是真实用户价值。\\
历史 live & M0/E-iter 72 请求、447812 tokens & 已发生的过程、成本和失败，不是新实验或盲测性能。\\
前瞻协议 & GP3/GP4 paired review 方案 & 应如何测量增量；尚没有测量结果。\\
\bottomrule
\end{tabularx}
\caption{事实层级。本文不会把某一层的成功转写成另一层尚不存在的证据。}
\label{tab:evidence_levels}
\end{table}

\subsection{版本、数据与时态}
实际 Git HEAD 为 \code{4dba622db459d4487bc8493b9af3a6a5f40a3387}；本报告描述其上的修后工作树，代码尚未提交。源码共 \ModuleCount{} 个 Python 文件、\SourceLines{} 行（含空行和注释），顶层测试文件 \TestFileCount{} 个，最终收集 \CollectedCases{} 个测试实例。文件数和行数是结构清点，不是能力覆盖率。当前完整 suite 在 414.20 秒通过；写报告不重跑模型或见证实验。\cite{acceptance,baseline}

2026-09-30 的 43 页报告记录修复前基线，用户随后提供的设计审查 PDF 是文档级诊断。此次源码核验将可复现问题、已有控制、产品功能缺口及条件任务逐项记录。\code{CURRENT_STATUS} 的旧段落按各自时间理解；最新状态以本文与 \code{audit/design_review_20260930/} 为准。

\subsection{格式参考与相关思想}
报告参考 OneReason 的单栏技术报告组织：紧凑摘要与总览图、显式设计理念、方法与协议、定量实验、部署/使用边界，以及可复现附录。\cite{onereason}图、数值、框图、伪代码和正文均根据 cv-agent 自身材料整理，不使用其品牌或实验结果。

Code-RAG 借用外部检索为生成提供可更新上下文的思想，但此项目的冻结 E2 默认是词法排序，不是该 RAG 论文中的稠密检索模型。\cite{rag}研究路径的交替模型动作与工具观察采用 ReAct 式组织；\cite{react}它增加了身份、必需验证完成、费用与语义准入要求。本文不把这些概念关系当作算法新颖性或效果增量证明。

\begin{insight}{读者可以直接依赖与仍需确认的部分}
可以依赖：默认静态路径的结构化产物、已枚举边界正反控制，以及候选绑定的有界事实。仍需确认：真实审查任务中的复核时间与事实质量、可选模型说明的净收益、私有 scratch/native 能力、跨仓库泛化和模型自主场景提出。原 G-A 为 NO；新的 GP0/GP1 仅通过静态工程门槛。
\end{insight}
